% titel: Blatt 8, A5: Lokale Extrema und Sattelpunkte
% Teil II: der durchgerechnete Weg.

\subsection*{Gegeben und gesucht}
a) Zeigen: $f$ hat bei $x_0$ einen Wendepunkt genau dann, wenn $f'$ bei
$x_0$ ein isoliertes lokales Extremum hat. b) Beweisen: aus
$f^{(n)}(x_0)=0$ für $n<m$ und $f^{(m)}(x_0)$ ungleich $0$ folgt -- $m$
gerade: isoliertes lokales Extremum, $m$ ungerade: isolierter
Sattelpunkt.

\subsection*{a) Wendepunkt und Extremum von $f'$}
\emph{Was ein Wendepunkt ist:} Die Kurve wechselt ihre Krümmung. Die
Krümmung wird durch $f''$ gesteuert; ein Krümmungswechsel bei $x_0$
bedeutet, dass $f''$ dort das Vorzeichen wechselt. Da $f''$ die Ableitung
von $f'$ ist, heisst genau: $f''$ wechselt das Vorzeichen bei $x_0$,
also: $f'$ hat bei $x_0$ ein \emph{lokales Extremum} (vorher steigt $f'$, nachher
fällt sie, oder umgekehrt).

\emph{Warum \emph{isoliert}:} Ein isoliertes lokales Extremum von $f'$ hat
eine Umgebung ohne weitere Extremalstellen. Das schliesst den Fall aus, dass
$f''\equiv0$ auf einem ganzen Intervall gilt (gerade Stücke, keine
Krümmungsänderung). Nur ein \emph{echter Wechsel} ist ein Wendepunkt. Damit
gilt die Äquivalenz: $f$ hat bei $x_0$ einen Wendepunkt genau dann, wenn
$f'$ bei $x_0$ ein isoliertes lokales Extremum hat.

\subsection*{b) Die Parität der ersten nicht verschwindenden Ableitung}
\emph{Kernargument (Taylor):} Wenn alle Ableitungen bis $m-1$ bei $x_0$
verschwinden und $f^{(m)}(x_0)$ ungleich $0$ ist, bildet $f^{(m)}(x_0)$
den ersten führenden Taylor-Term. In der Nähe von $x_0$ gilt
\begin{equation*}
  f(x)-f(x_0)
  =\frac{f^{(m)}(x_0)}{m!}\,(x-x_0)^m+\text{kleinere Reste},
\end{equation*}
und für genügend kleine Abstände beherrscht der Term $(x-x_0)^m$ das
Vorzeichen von $f(x)-f(x_0)$.

\begin{itemize}
  \item \textbf{$m$ gerade:} $(x-x_0)^m\ge0$ links \emph{und} rechts von
        $x_0$. Das Vorzeichen von $f(x)-f(x_0)$ ist damit auf beiden
        Seiten dasselbe (das von $f^{(m)}(x_0)$): ein \emph{isoliertes
        lokales Extremum}. Minimum, falls $f^{(m)}>0$; Maximum, falls
        $f^{(m)}<0$.
  \item \textbf{$m$ ungerade:} $(x-x_0)^m$ wechselt das Vorzeichen bei
        $x_0$. Die Differenz $f(x)-f(x_0)$ hat links und rechts
        \emph{entgegengesetzte} Vorzeichen: kein Extremum. Da $m\ge3$
        (wegen $f'(x_0)=0$), ist $x_0$ mit waagerechter Tangente ein
        \emph{isolierter Sattelpunkt}.
\end{itemize}

\emph{Induktive Sicht (warum man auf Ordnung $m-1$ reduzieren kann):}
Dasselbe Ergebnis lässt sich induktiv führen, indem man die Paritätsaussage
auf die Ableitung $g=f'$ anwendet: Gelten für $f$ die Bedingungen zur
Ordnung $m+1$ (erste $m$ Ableitungen verschwinden, $f^{(m+1)}$ ungleich $0$), so
gelten für $g=f'$ die Bedingungen zur Ordnung $m$ (denn $g^{(n)}=f^{(n+1)}$
und $g^{(m)}=f^{(m+1)}$ ungleich $0$). Die führende Ordnung sinkt also pro
Ableitung um $1$, und damit kehrt sich die Parität gerade/ungerade um:
Aus einem lokal extremalen $f'$ ($m$ gerade) wird ein Wendepunkt-artiges
$f$ ($m+1$ ungerade), aus einem sattel-artigen $f'$ ($m$ ungerade) ein
extremales $f$ ($m+1$ gerade). Das liefert im Kern dieselbe Aussage wie
die Taylor-Rechnung und ist gerade der Weg in der Vorlesung: Parität der
führenden Ordnung von $f'$ entscheidet, ob $f'$ ein Extremum oder einen
Durchgang hat -- und damit, ob $f$ ein Extremum oder einen Sattelpunkt
besitzt.

\textit{Abschluss des Beweises:} Die Taylor-Rechnung oben zeigt beide
Richtungen der Paritätsaussage vollständig; die induktive Sicht ist
verständnismässig dasselbe Argument, nur über die Ableitung (Ordnung $m-1$)
geführt.

\begin{ergebnis}
\begin{equation*}
  f\ \text{hat bei }x_0\ \text{einen Wendepunkt} \iff
  f'\ \text{hat bei }x_0\ \text{ein isoliertes lokales Extremum}.
\end{equation*}
Aus $f^{(n)}(x_0)=0\\ (n<m)$ und $f^{(m)}(x_0)$ ungleich $0$: für gerades $m$ ein
isoliertes lokales Extremum, für ungerades $m$ ein isolierter Sattelpunkt.
\end{ergebnis}

\begin{kontrolle}
\textbf{Beispiel} $f(x)=(x-x_0)^3$: $f'=3(x-x_0)^2$ hat bei $x_0$ ein
Minimum, $f''(x_0)=0$, $f'''(x_0)=6$ ungleich $0$, $m=3$ ungerade --
Sattelpunkt, wie behauptet. \checkmark

\textbf{Beispiel} $f(x)=(x-x_0)^4$: $f'''(x_0)=0$, $f^{(4)}(x_0)=24$ ungleich $0$,
$m=4$ gerade -- isoliertes lokales Minimum. \checkmark

\textbf{Grenzfall:} Wäre $f^{(m)}(x_0)=0$ für \emph{alle} $m$, so gäbe es
keine erste nicht verschwindende Ableitung und keine Vorzeichenaussage;
die Zusatzannahme $f^{(m)}(x_0)$ ungleich $0$ ist unverzichtbar. \checkmark
\end{kontrolle}

\begin{merksatz}
Der Wendepunkt ist ein isoliertes lokales Extremum von $f'$. Entscheidend
ist die \emph{Parität} der ersten nicht verschwindenden Ableitung: gerade
Ordnung gibt ein Extremum (Term $(x-x_0)^m$ ohne Vorzeichenwechsel),
ungerade einen Sattelpunkt (Term mit Vorzeichenwechsel). Über $g=f'$ lässt
sich die Aussage induktiv auf Ordnung $m-1$ zurückführen.
\end{merksatz}
